%======================================================================
%  Discussion
%
%  NEW 2026-09-04. No companion artifact supplies Discussion prose, so this
%  section was assembled from material that already existed in the manuscript
%  and was sitting in the wrong place. Provenance of each paragraph:
%
%   P1  summary of contribution -- restates Results, no new claim.
%   P2  ensemble mechanism. Moved out of old 03_origin.tex 3.4 P2. The author
%       corrected the framing on 2026-09-04: that bacterial chromosomes carry
%       integrated plasmids and other MGEs is an ESTABLISHED FINDING, not a
%       hypothesis, and must be cited as such. The hypothesis is the narrower
%       one that Evo2 learned signal from well-sampled bacterial genomes that
%       IMG/PR does not contain. Rewritten accordingly.
%   P3  interpretation limitations. Consolidated from six passages that were
%       scattered through 05_interpretation.tex.
%   P4  benchmark and deployment limitations. The first three are the
%       "Limitations (state proactively)" bullets from the origin artifact's
%       "Open items" band. CONVERSION_NOTES.md flagged them as reading closer
%       to manuscript material than to memo and said they "may deserve a home
%       in the Discussion". This is that home. Author sign-off requested.
%   P5  outlook.
%
%  Nature Biotechnology does not use subheadings in the Discussion.
%======================================================================

\section*{Discussion}
\label{sec:discussion}

\model shows that a genomic language model trained only on plasmids, and
sampled for the breadth of plasmid diversity rather than for sequencing
depth, learns representations from which the non-coding elements governing
plasmid transfer and replication can be read out by a lightweight probe. The
practical consequence is a deployable annotation route rather than a
benchmark score: the same two-stage pipeline that we evaluate on curated
origins runs unchanged across IMG/PR, reduces each plasmid to a short ranked
list of candidate regions, and does so at a cost that makes exhaustive
application to a metagenomic database tractable.

Combining \model with Evo2-7B improved on either model alone, which indicates
that the two carry partly non-redundant predictive information. Two
differences in pretraining corpus could produce that. \model was pretrained on
plasmids with diversity-aware sampling that preserves within-PTU variation,
whereas Evo's published OpenGenome construction retained a single IMG/PR
representative per PTU, so the two models saw different amounts of
within-clade structure. Evo2 was also pretrained on whole prokaryotic genomes,
and bacterial chromosomes are well documented to carry integrated plasmids,
integrative and conjugative elements, and other mobile genetic elements
% The load-bearing citations here are Guglielmini (ICEs are the most abundant
% conjugative elements in prokaryotes, from >1,000 genomes), Cury (integrated
% and extrachromosomal elements interconvert) and Johnson (ICEs encode their
% own oriT and relaxase). The author-supplied Yan preprint is a prevalence
% figure (25-75% of single cells carry an MGE), so it is cited last as support
% rather than as the primary claim. It has no peer-reviewed version yet.
\cite{guglielmini2011repertoire,cury2018host,johnson2015integrative,yan2025revealing},
so a chromosomal corpus is not free of origin-like and transfer-associated
sequence. Our hypothesis is therefore the narrower one that Evo2 acquired
signal from well-sampled bacterial genomes that the IMG/PR plasmid corpus does
not contain. Testing it would mean stratifying the ensemble gain by the
presence of related chromosomally integrated elements, or by other proxies for
corpus-specific coverage.

The interpretation analyses should be read as illustrative rather than as a
systematic survey. The loci shown were selected because their architecture is
documented, no quantitative test was run across a defined denominator of loci,
and each dependency map carries its own local colour scale, so pattern
intensities cannot be compared between panels. Two case studies do not
establish that the correspondence holds for origins generally, the four
promoters were chosen as visually clean examples from a motif scan and do not
support a promoter-wide recovery rate, and no element-recovery rate is claimed
anywhere. The same applies to the embedding projection: no separation metric
or classifier was computed, the cohort was selected for curated hits and is
not a representative sample of plasmids, and because the projection varies
strongly with GC content, sequence composition may contribute to the apparent
annotation structure.

Three limitations bound the origin results themselves. Circular plasmids are
linearized for windowing, so a small number of origins spanning the
linearization point are represented differently from the way the replicating
molecule presents them. The curated origin sets are drawn from databases with
partly overlapping content, so redundancy between them is controlled by
filtering and by similarity-aware partitioning rather than eliminated. And the
pipeline ranks origin-containing regions rather than delimiting origins: a
top-ranked region localizes an origin to an intergenic interval, not to its
boundaries, and downstream use should treat a prediction as a candidate for
experimental confirmation. On \emph{oriV} in particular, the margin over a
baseline that ranks the longest intergenic region first is real but modest,
and a reference-based tool remains competitive.

The elements this model annotates are the two that determine whether a plasmid
can be moved between hosts and maintained in one, which is precisely what
makes an uncharacterized plasmid usable as an engineering vector. Applied at
database scale to sequences whose hosts have never been cultured, an
alignment-free route to those elements turns the metagenomic plasmid record
from a catalogue of sequences into a searchable source of candidate parts.
